\begin{helper}
Let \(G\) be a loop graph, and let \(T_G\) be the Euclidean adjacency operator
from \(\noderef{LoopGraphAdjacencyEuclideanOperator}\).  If \(x\) is a nonzero
vector in the zero-sum subspace \(W_0(V(G))\) and
\[
  T_Gx=\mu x,
\]
then \(\mu\) is a non-principal eigenvalue of \(G\) in the sense of
\(\noderef{LoopGraphNonprincipalEigenvalue}\).
\end{helper}
\begin{proof}
Define \(f\colon V(G)\to\mathbb R\) by \(f(v)=x(v)\).  Since \(x\neq 0\) in
the zero-sum subspace, some coordinate of \(x\) is nonzero; otherwise every
coordinate would be zero and function extensionality would give \(x=0\).  Thus
there is a vertex \(v\) with \(f(v)\neq 0\).

The membership \(x\in W_0(V(G))\), using
\(\noderef{LoopGraphEuclideanZeroSumSubmodule}\), gives
\[
  \sum_{v\in V(G)} f(v)=0.
\]
The eigen-equation \(T_Gx=\mu x\) gives, coordinate by coordinate,
\[
  (T_Gx)(v)=\mu x(v).
\]
By the definition of \(T_G\) in \(\noderef{LoopGraphAdjacencyEuclideanOperator}\),
the left side is the adjacency action of \(G\) applied to the coordinate
function \(f\).  Therefore
\[
  \text{adjacency-action}_G(f)(v)=\mu f(v)
\]
for every \(v\).  The function \(f\) is nonzero, has zero coordinate sum, and
satisfies the adjacency eigenvector equation, which is exactly the definition
of \(\noderef{LoopGraphNonprincipalEigenvalue}\).
\end{proof}
